\documentclass[journal,10pt,twocolumn]{IEEEtran}
\usepackage{cite}
\usepackage{amsmath,amssymb,amsfonts}
\usepackage{algorithmic}
\usepackage{graphicx}
\usepackage{textcomp}
\usepackage{xcolor}
\usepackage{booktabs}
\usepackage{multirow}
\usepackage{hyperref}

\begin{document}

% 1) TITLE
\title{ALAS-SCLF: Adaptive-L Asymmetric-Segmented SCLF Polar Decoders with Guard-Band Stabilization}

% 2) AUTHOR
\author{Deep~Shekhar~Halder,~\IEEEmembership{Adamas University,}
        and~Dwaipayan~Ghosh,~\IEEEmembership{Kalyani Government Engineering College}%
\thanks{Deep Shekhar Halder is with the Department of Computer Science and Engineering, Adamas University, Kolkata, West Bengal, India (e-mail: deep.halder@adamasuniversity.ac.in).}%
\thanks{Dwaipayan Ghosh is with the Department of Electronics and Communication Engineering, Kalyani Government Engineering College, Kalyani, West Bengal, India (e-mail: dwaipayan.ghosh@kgec.ac.in).}}

\maketitle

% 3) ABSTRACT
\begin{abstract}
This paper presents a complete technical exposition of the \textbf{ALAS-SCLF (Adaptive-List Asymmetric-Segmented Successive Cancellation List Flip)} decoding framework for polar codes in 5G New Radio (NR) Ultra-Reliable Low-Latency Communication (URLLC) networks. While threshold-learning-based SCLF (TS-SCLF) decoders integrate intermediate segmented cyclic redundancy checks (CRCs) and early termination (SCL-RE) to reduce decoding latency, they exhibit critical limitations: (1) post-flip path-metric transient instability, (2) search-space lockouts caused by rigid symmetric CRC partitioning, and (3) rigid search list constraints under high noise. ALAS-SCLF resolves these bottlenecks by introducing: (i) a post-flip stabilization Guard-Band ($G = \lceil \beta \cdot N \rceil$), (ii) asymmetric segmented CRC partitioning ($l_{\text{seg1}} = \lceil \gamma \cdot K \rceil$) with Segment-Constrained Flipping (SCF), (iii) adaptive list-size scaling ($L = 4 \rightarrow 8$), and (iv) Deep Q-Network (DQN) parameter optimization. Extensive Monte Carlo simulations conducted under a preferred statistical baseline of 20,000 trials per SNR point across short block $(128, 56)$ and long block $(1024, 512)$ examples demonstrate that ALAS-SCLF achieves up to a 36.6\% relative reduction in Frame Error Rate (FER) at 1.5 dB SNR for $(128, 56)$, a 43.2\% FER reduction for $(1024, 512)$, and up to a 34.3\% reduction in average decoding attempts over TS-SCLF, establishing a state-of-the-art Pareto frontier.
\end{abstract}

% 4) KEYWORD
\begin{IEEEkeywords}
Polar Codes, Successive Cancellation List (SCL), Bit-Flipping, Asymmetric CRC, Guard-Band, Deep Q-Network (DQN), 5G NR URLLC, Monte Carlo 20000 Trials.
\end{IEEEkeywords}

% 5) INTRODUCTION
\section{Introduction}
\IEEEPARstart{I}{n} modern wireless communication standards, such as 5G New Radio (NR), the requirements for Ultra-Reliable Low-Latency Communications (URLLC) mandate highly efficient channel coding schemes. Polar codes, introduced by Arikan, have been adopted as the channel coding scheme for control channels in 5G NR due to their capacity-achieving performance under successive cancellation (SC) decoding. To improve finite block-length performance, successive cancellation list (SCL) decoding is combined with a cyclic redundancy check (CRC) helper.

To bridge the performance gap between low-complexity SC and high-performance SCL with large list sizes, Successive Cancellation List Flip (SCLF) decoders restart SCL decoding by flipping error-prone bits when the initial SCL run fails the CRC. Recently, a Threshold-Learning-Based SCLF (TS-SCLF) decoder was proposed to minimize decoding latency by introducing intermediate segmented CRCs and early SCL termination. However, TS-SCLF suffers from three main limitations:
\begin{enumerate}
    \item \textbf{Post-Flip Metric Instability:} Right after a bit is forcefully flipped at layer $l_{\text{flip}}$, the path metrics (PMs) of candidate paths undergo a severe transient period. Checking early termination thresholds immediately causes the correct path to be prematurely pruned because the path metric spread has not stabilized.
    \item \textbf{False-Pass Search Space Lockout:} TS-SCLF uses a symmetric 50/50 partition of the information bits. A longer segment size leads to high intermediate CRC false-pass probability under noise. If Segment 1 falsely passes the check despite errors, the decoder is locked out of search-space flips inside Segment 1, leading to a decoding failure.
    \item \textbf{Rigid Search List Constraints:} Under severe channel noise, keeping the list size restricted to $L = 4$ restricts search capacity, causing the decoder to fail to identify the correct path even with multiple flips.
\end{enumerate}

\section{Objective}
The primary objective of this internship research project is to design, implement, and evaluate the Adaptive-List Asymmetric-Segmented SCLF (ALAS-SCLF) decoder framework to resolve the core weaknesses of threshold-learning segmented flip decoders in 5G NR control channels:
\begin{enumerate}
    \item Mitigate the post-flip path metric spread transient instability through a stabilization Guard-Band ($G$) window.
    \item Prevent search-space lockouts by introducing asymmetric CRC partitioning and Segment-Constrained Flipping (SCF).
    \item Enhance error correction capability under severe channel noise scenarios by dynamically scaling SCL list sizes.
    \item Integrate a Deep Q-Network (DQN) reinforcement learning agent to automate parameter search and optimization.
\end{enumerate}

% 6) LITERATURE / BACKGROUND
\section{Literature / Background}
This section establishes the mathematical foundations of polar coding, successive cancellation decoding, SCL list extension, and the Deep Q-Network (DQN) optimization framework. All formulations are generic and applicable to any codeword length $N = 2^n$ and information bits $K$.

\subsection{Polar Codes}
Polar codes polarize $N$ independent subchannels into capacity-achieving subchannels. The polar generator matrix $\mathbf{G}_N$ is constructed via the Kronecker power of the kernel matrix $\mathbf{F}$:
\begin{equation}
\mathbf{G}_N = \mathbf{B}_N \mathbf{F}^{\otimes n}, \quad \mathbf{F} = \begin{bmatrix} 1 & 0 \\ 1 & 1 \end{bmatrix}
\end{equation}
where $\mathbf{B}_N$ is the $N \times N$ bit-reversal permutation matrix. The $K$ information payload bits are mapped to the high-reliability subchannels index set $\mathcal{I}$. The remaining $N - K$ subchannels represent the frozen set $\mathcal{I}^c$, fixed to zero. The transmitted codeword vector is computed as $\mathbf{x} = \mathbf{u} \mathbf{G}_N$.

\subsection{SC Based decoding}
In successive cancellation, decoding is performed recursively over the binary decoding tree. Let $\mathbf{P} = (P_1, P_2, \dots, P_M)$ represent the LLR vector received from a parent node of length $M$. The left child node LLRs ($P_{\text{left}, i}$) and right child node LLRs ($P_{\text{right}, i}$) are recursively updated for $i \in \{1, 2, \dots, M/2\}$ as:
\begin{equation}
P_{\text{left}, i} = 2 \operatorname{arctanh}\left( \tanh\frac{P_{2i-1}}{2} \cdot \tanh\frac{P_{2i}}{2} \right)
\end{equation}
\begin{equation}
P_{\text{right}, i} = P_{2i} + (1 - 2\hat{u}_{\text{left}, i}) \cdot P_{2i-1}
\end{equation}
where $\hat{u}_{\text{left}, i}$ represents the hard-decision feedback bit computed from the left child node. At each leaf node $l$ corresponding to an information bit ($l \in \mathcal{I}$), the active path splits into two candidate paths representing $u_l = 0$ and $u_l = 1$. The path metric PM for path $m$ is updated recursively according to:
\begin{equation}
\text{PM}_l^{(m)} = \text{PM}_{l-1}^{(m)} + \begin{cases} 0, & \text{if } u_l^{(m)} = \operatorname{Hard}(P_l^{(m)}) \\ |P_l^{(m)}|, & \text{otherwise} \end{cases}
\end{equation}

\subsection{DeepQ-Network (DQN)}
To automatically discover optimal parameters, we formulate a Markov Decision Process (MDP) defined by the tuple $(\mathbf{S}, \mathbf{A}, P, R, \gamma_{\text{RL}})$:
\begin{enumerate}
    \item \textbf{State Space (S):} The state vector represents the parameter configuration $\mathbf{s}_t = [D_{1, t}, D_{2, t}, \alpha_t]$.
    \item \textbf{Action Space (A):} The action space consists of discrete adjustments $\mathbf{a}_t \in \{ \pm\Delta D_1, \pm\Delta D_2, \pm\Delta\alpha \}$.
    \item \textbf{Reward Function (R):} Evaluates error-correction performance and attempts compared to reference values:
    \begin{equation}
    R(\mathbf{s}) = \begin{cases} T_{\text{prev}} - T_{\text{curr}}(\mathbf{s}), & \text{if } \text{FER}(\mathbf{s}) \le (1 + \delta) \text{FER}_{\text{ref}} \\ -1000, & \text{otherwise} \end{cases}
    \end{equation}
    \item \textbf{Target Q-Value Update:} Trained by minimizing the Bellman error using weights $\theta$:
    \begin{equation}
    Y_t = R_t + \gamma_{\text{RL}} \cdot \max_{a'} Q_{\text{target}}(\mathbf{s}_{t+1}, a'; \theta^-)
    \end{equation}
\end{enumerate}

% 7) PROPOSED ALAS-SCLF (GENERIC CONCEPT)
\section{Methodology}
This section describes the generic, value-independent concepts of the proposed ALAS-SCLF decoder. There are no specific numerical values assumed here; it represents the generic framework applicable to any polar code configuration $(N, K)$.

\subsection{Understanding of Work}
ALAS-SCLF addresses the metric stabilization and false-pass lockout issues of segmented decoders via three key concepts:
\begin{enumerate}
    \item \textbf{Post-Flip Guard-Band ($G$):} Flipped bits cause a transient non-stationary shift in LLR and path metrics. Checking thresholds immediately at layer $l_{\text{flip}} + 1$ prunes correct paths. We disable early termination threshold checks for a Guard-Band window $G$ defined as:
    \begin{equation}
    \text{Early Term Active iff: } l > l_{\text{flip}} + G, \quad G = \lceil \beta \cdot N \rceil
    \end{equation}
    where $\beta \approx 0.06$ is the channel stabilization coefficient.
    
    \item \textbf{Asymmetric Partitioning \& Segment-Constrained Flipping (SCF):} The payload $K$ is partitioned asymmetrically into Segment 1 of length $l_{\text{seg1}} = \lceil \gamma \cdot K \rceil$ ($\gamma \approx 0.28$) protected by intermediate CRC $c_1$, and Segment 2 protected by outer CRC $c_2$. If Segment 1 CRC fails, the candidate flipping set $S$ is strictly constrained to Segment 1:
    \begin{equation}
    S \subseteq [1, l_{\text{seg1}}], \quad \text{if } \text{CRC}_{\text{seg1}} = \text{Fail}
    \end{equation}
    
    \item \textbf{Adaptive List Scaling:} Channel noise severity is quantified after the initial decoding run using path metric range $R = \text{PM}_{\text{max}}(0) - \text{PM}_{\text{min}}(0)$. The list size for all restart runs is dynamically assigned as:
    \begin{equation}
    L_{\text{restart}} = \begin{cases} 8, & \text{if } R < R_{\text{thresh}} \text{ OR } \text{CRC}_{\text{seg1}} = \text{Fail} \\ 4, & \text{otherwise} \end{cases}
    \end{equation}
\end{enumerate}

\subsection{Algorithm}
ALAS-SCLF utilizes Algorithm 1 (FLIPSKIP) for path-flipping skip criterion, Algorithm 2 (SCL-RE-GB) for early termination with Guard-Band, and Algorithm 3 (Generic ALAS-SCLF Framework) for orchestrating the overall decoding process.

\subsection{System diagram}
\begin{figure}[h]
\centering
\includegraphics[width=\linewidth]{figures/fig_system_flowchart.png}
\caption{Detailed system architecture and decision flowchart for the proposed generic ALAS-SCLF decoder.}
\end{figure}

% 8) RESULTS & DISCUSSION
\section{Results \& Discussion}
The performance of the proposed ALAS-SCLF decoder was evaluated using extensive code-driven Monte Carlo simulations over an AWGN channel under BPSK modulation. To establish maximum statistical confidence, all evaluations enforce a preferred sample baseline of 20,000 Monte Carlo trials per SNR point across multiple specific code parameters.

\subsection{Short Block Polar Code: $(128, 56)$}
At SNR = 1.5 dB, baseline TS-SCLF exhibits a severe FER degradation (crossover at 0.02050 compared to 0.01650 for SCLF). This occurs because symmetric $l_{\text{seg}} = 24$ leads to false Segment 1 CRC passes, locking the decoder into searching Segment 2 when the error resides in Segment 1. ALAS-SCLF resolves this, dropping FER to 0.01300 (a 36.6\% relative FER reduction over TS-SCLF) and cutting attempts at 0.5 dB from 2.162 to 1.707 (a 21.0\% latency reduction).

\subsection{Medium/Long Block Polar Code: $(1024, 512)$}
For medium/long block lengths typical of 5G data channels, ALAS-SCLF achieves even greater error-rate gains. At SNR = 1.5 dB, ALAS-SCLF reduces FER from 0.01250 to 0.00710 (a 43.2\% relative FER reduction). Average decoding attempts at SNR = 1.0 dB drop from 2.890 down to 1.900 (a 34.3\% complexity reduction), proving superior scalability across block sizes.

% 9) CONCLUSION
\section{Conclusion}
This paper presented the generic theoretical formulation, visual operational flowchart, algorithmic specifications, and multi-example simulation evaluation of ALAS-SCLF. By integrating post-flip Guard-Band stabilization ($G = \lceil \beta \cdot N \rceil$), asymmetric CRC partitioning ($l_{\text{seg1}} = \lceil \gamma \cdot K \rceil$), and adaptive list scaling ($L=4 \rightarrow 8$), ALAS-SCLF eliminates false-pass lockouts and premature path pruning, establishing a new state-of-the-art Pareto frontier for 5G URLLC decoders.

% 10) REFERENCE
\begin{thebibliography}{10}
\bibitem{arikan2009} E. Arikan, "Channel polarization: A method for constructing capacity-achieving codes for symmetric binary-input memoryless channels," \emph{IEEE Trans. Inf. Theory}, vol. 55, no. 7, pp. 3051-3073, Jul. 2009.
\bibitem{tal2015} I. Tal and A. Vardy, "List decoding of polar codes," \emph{IEEE Trans. Inf. Theory}, vol. 61, no. 5, pp. 2213-2226, May 2015.
\bibitem{afisiadis2014} O. Afisiadis et al., "A low-complexity improved successive cancellation decoder for polar codes," in \emph{Proc. Asilomar Conf. Signals, Syst. Comput.}, Nov. 2014, pp. 2116-2120.
\bibitem{chandesris2016} L. Chandesris, V. Savin, and D. Declercq, "An improved SC-flip decoder for polar codes," \emph{IEEE Commun. Lett.}, vol. 20, no. 12, pp. 2333-2336, Dec. 2016.
\bibitem{sun2026} Y. Sun and C. Zhang, "TS-SCLF: Threshold-Learning-Based SCLF Polar Decoder With Segmented CRC," \emph{IEEE Wireless Commun. Lett.}, vol. 15, no. 5, pp. 681-684, May 2026.
\bibitem{zhou2016} H. Zhou et al., "Segmented CRC-aided SC list polar decoding," in \emph{Proc. IEEE Veh. Technol. Conf.}, Jun. 2016, pp. 1-5.
\bibitem{ueng2018} Y. L. Ueng et al., "Successive cancellation list bit-flip decoder for polar codes," in \emph{Proc. IEEE WCSP}, Oct. 2018, pp. 1-6.
\bibitem{liu2024} Y. J. and R. Liu, "Reducing complexity of SC-based flip decoding of polar codes by early-stopping," \emph{IEEE Commun. Lett.}, vol. 28, no. 4, pp. 768-772, Apr. 2024.
\bibitem{liang2024} F.-S. Liang et al., "Deep-learning-aided successive cancellation list flip decoding for polar codes," \emph{IEEE Trans. Cogn. Commun. Netw.}, vol. 10, no. 2, pp. 374-386, Apr. 2024.
\bibitem{li2024} J. Li et al., "Deep learning-assisted adaptive dynamic-SCLF decoding for polar codes," \emph{IEEE Trans. Cogn. Commun. Netw.}, vol. 10, no. 3, pp. 836-851, Jun. 2024.
\end{thebibliography}

\end{document}
